\section{Path With Minimum Effort}
\label{sec:graph-path-min-effort}

\problemstatement{Given a 2D grid of heights, find a path from the
top-left cell to the bottom-right cell (moving $4$-directionally)
that minimizes the \textbf{effort} of the path, where a path's effort
is defined as the \textbf{maximum} absolute height difference between
any two consecutive cells along it.}

\begin{patternbox}
\emph{``Minimize the maximum edge weight along a path''} (a
\textbf{minimax path}) is solved with the exact same frontier-expansion
machinery as Section~\ref{sec:graph-dijkstra}'s Dijkstra, with one
change: the quantity tracked and relaxed is no longer a
\textbf{sum} of edge weights, but the \textbf{running maximum} edge
weight seen so far along the path.
\end{patternbox}

\subsection*{Adapting the relaxation rule}

Let $\textit{effort}[\textit{cell}]$ be the minimum possible
``maximum height difference'' achievable on some path from the start
to $\textit{cell}$. Moving from $u$ to a neighbor $v$, the effort of
extending that path is $\max(\textit{effort}[u],
|\textit{height}[v]-\textit{height}[u]|)$ -- not
$\textit{effort}[u] + |\textit{height}[v]-\textit{height}[u]|$, since a
path's effort is its single worst step, not an accumulated total. The
relaxation check becomes: if this candidate effort is
\textbf{smaller} than the currently known $\textit{effort}[v]$, update
and push.

\subsection*{Why Dijkstra's greedy-frontier property still applies}

The same correctness argument from Section~\ref{sec:graph-dijkstra}
carries over unchanged: popping the cell with the smallest known
effort from the priority queue guarantees that effort is final,
because every alternative path to it would have to pass through a
cell with effort $\ge$ the popped value, and taking the
$\max$ with any additional step can only keep the running effort the
same or increase it -- never decrease it. This ``can only stay the
same or grow'' property is the minimax analogue of Dijkstra's
``non-negative weights only add'' requirement.

\subsection*{C++ implementation}

\begin{lstlisting}
#include <bits/stdc++.h>
using namespace std;

int minimumEffortPath(vector<vector<int>>& heights) {
    int rows = heights.size(), cols = heights[0].size();
    vector<vector<int>> effort(rows, vector<int>(cols, INT_MAX));
    effort[0][0] = 0;

    priority_queue<tuple<int,int,int>, vector<tuple<int,int,int>>, greater<>> pq;
    pq.push({0, 0, 0}); // {effort, row, col}

    int dr[] = {-1, 1, 0, 0};
    int dc[] = {0, 0, -1, 1};

    while (!pq.empty()) {
        auto [e, r, c] = pq.top();
        pq.pop();
        if (e > effort[r][c]) continue;
        if (r == rows - 1 && c == cols - 1) return e;

        for (int dir = 0; dir < 4; dir++) {
            int nr = r + dr[dir], nc = c + dc[dir];
            if (nr >= 0 && nr < rows && nc >= 0 && nc < cols) {
                int candidate = max(e, abs(heights[nr][nc] - heights[r][c]));
                if (candidate < effort[nr][nc]) {
                    effort[nr][nc] = candidate;
                    pq.push({candidate, nr, nc});
                }
            }
        }
    }
    return 0;
}

int main() {
    vector<vector<int>> heights = {
        {1, 2, 2},
        {3, 8, 2},
        {5, 3, 5}
    };
    cout << minimumEffortPath(heights) << endl; // 2
    return 0;
}
\end{lstlisting}

\subsection*{Dry run}

\dryrun{Take the $3\times3$ heights grid above.}

The path $(0,0)\to(0,1)\to(0,2)\to(1,2)\to(2,2)$ has consecutive
differences $|2{-}1|=1$, $|2{-}2|=0$, $|2{-}2|=0$, $|5{-}2|=3$ --
effort $3$. A better path,
$(0,0)\to(1,0)\to(2,0)\to(2,1)\to(2,2)$, has differences
$|3{-}1|=2$, $|5{-}3|=2$, $|3{-}5|=2$, $|5{-}3|=2$ -- effort $2$,
worse on individual steps but capped lower overall. Dijkstra's
frontier expansion explores cells in increasing order of their best
known effort, popping $(2,2)$ (the destination) only once its true
minimum effort, $2$, has been finalized.

\begin{complexitybox}
\textbf{Time Complexity.} $O(\textit{rows} \cdot \textit{cols} \cdot
\log(\textit{rows}\cdot\textit{cols}))$, identical structure to
Dijkstra over a grid graph with $4$ edges per node.
\textbf{Space Complexity.} $O(\textit{rows}\cdot\textit{cols})$ for
the \textit{effort} grid and the priority queue.
\end{complexitybox}

\begin{takeawaybox}
Dijkstra's algorithm generalizes to any relaxation rule that is
\textbf{monotonic} -- extending a path by one more edge can only keep
the tracked quantity the same or make it worse, never better. Sum of
weights and running maximum both satisfy this, which is why the exact
same priority-queue frontier expansion solves both the standard
shortest-path problem and this minimax variant.
\end{takeawaybox}
